% Data and Software Availability, and the disclosure of language-model use.
% Labels defined here: sec:avail, sec:llm.
% Every locator comes from artifact.tex (\fsArtifactLink, \fsUpstreamZenodo);
% nothing is typed in by hand.

\section*{Data and Software Availability}
\label{sec:avail}

The complete artifact for this work is deposited at \fsArtifactLink. It holds
the source code, the environment lockfile with the exact version of every
dependency, the scripts that fetch the input data, the sha256 hash of every
dataset file and checkpoint, the split manifests, the probe seeds, every raw
per-structure prediction and every timing record, the scripts that draw each
figure and table, the test suite, and a worked example. The code is released
under the MIT licence and the records under CC~BY~4.0. File paths below are
given relative to the root of the artifact.

The inputs are not ours and are not redistributed. The committee checkpoints
and the 3BPA, rMD17 and liquid-water frames are the supporting data of
\citet{beck2025committee}, archived at \texttt{\fsUpstreamZenodo}. The MPtraj
structures are the ones distributed with CHGNet~\citep{deng2023chgnet}. The
OMat24 validation structures used in the out-of-distribution scan are from
\citet{barrosoluque2024omat24} and are released under CC~BY~4.0. The fetch
scripts download each file and compare it with its recorded hash before any
use, so a substituted or corrupted input stops the run instead of changing a
result silently.

Every statistical result regenerates on a CPU. The command
\texttt{tools/reproduce.py --tier offline} re-derives each headline quantity
from the committed records, with no model execution, no CUDA and no download,
and checks each quantity against its committed value. \texttt{tools/audit.py}
verifies the frozen prior-work tree, the sha256 binding of each pre-registered
protocol to a registration record written before that experiment produced any
output, the provenance of every number in this manuscript, the split manifests
and the test-access ledger. \texttt{tools/report\_go\_no\_go.py} prints every
scored prediction as data, next to the confidence that was stated for it in
advance.

Three kinds of computation need a GPU. Regenerating the cached per-head force
matrices for the molecular and water systems needs one GPU; everything
downstream of that cache is linear algebra on the CPU. The materials scans
need one GPU each: the MPtraj scan ran on one A100, and the OMat24 scan and
the active-learning round on 3BPA ran on partitioned (MIG) slices of A100
cards. Such slices are suitable for predictions and for training but not for
timing, so no wall-clock number in this paper comes from a slice. The timing
experiments of Section~\ref{sec:wallclock} ran on A100-SXM4-80GB, H200 and
RTX~PRO~6000 Blackwell cards under one pinned toolchain. The device, driver,
streaming-multiprocessor count and clock state of every timing run are
recorded next to the measurements of that run, because the wall-clock
conclusions are properties of the hardware as much as of the method.

The records added for this revision are included in the artifact. They are
the per-structure predictions, reference forces and scores of the
out-of-distribution scan on OMat24 (Table~\ref{tab:ood}); the selected sets,
retraining logs and held-out errors of the active-learning round on 3BPA
(Table~\ref{tab:al}); and the accuracy and calibration record
\texttt{results/records/j12\_accuracy.json} behind Table~\ref{tab:accuracy}
and Figure~\ref{fig:calibration}. The protocols for the two new experiments
were registered before either experiment produced any output, and the
registration records are included alongside the results.

\section*{Use of Large Language Models}
\label{sec:llm}

Large-language-model coding agents were used substantially throughout this
work and throughout this revision: for research ideation, for writing the
experiment scripts and the analysis code, for the adversarial review described
below, and for drafting and revising the manuscript. The author directed the
work, chose the questions, and is responsible for the contents.

Because that exposure is broad, verification is mechanical rather than a
matter of inspection. Every number in this paper is generated from a committed
record by \texttt{tools/paper\_numbers.py} and re-checked by
\texttt{tools/audit.py}; none is transcribed, so an agent cannot introduce a
value by typing one. The protocol of each experiment was registered and
hash-bound to a record written before that experiment produced any output, so
a hypothesis cannot be adjusted after the outcome is known. Split roles are
content-hashed and recompute from recorded seeds. An adversarial review,
itself agent-driven and documented in \texttt{results/REVIEW\_FINDINGS.md},
found real defects in agent-written code and in the supporting documents,
among them a benchmark that measured fewer backward passes than it reported
and a filename collision that destroyed a result file; a cross-document audit
found many stale numbers in the supporting documents. The macro-derived
numbers, which are the ones in this manuscript, were not among the failures.
That record is the reason every number here is derived rather than copied.
